\section{Analytic regularity of the local solution map}
\label{sec:regularity}
The deformation estimates control how far an exact root moves under a small
perturbation. We now show that its local dependence on the parameters is
analytic, and compute the derivative from the polynomial Jacobian. The
analytic argument takes place over a complete normed field. For integral
$p$-adic data, we first include the coefficients and base point in $\Qp$.
$\Zp$ itself is not a vector space over $\Qp$.

Let $\mathbb K$ be a complete nontrivially normed field. Let $e,k\in\N$, and
write the variables as $(Y,T)\in\mathbb K^e\times\mathbb K^k$. For
\[
 F\in\mathbb K[Y_1,\ldots,Y_e,T_1,\ldots,T_k]^e,
 \qquad a=(y_0,t_0),
\]
define the matrices of evaluated formal partial derivatives
\begin{equation}
 J=D_YF(a)\in M_e(\mathbb K),\qquad
 K=D_TF(a)\in M_{e\times k}(\mathbb K).
 \label{eq:regularity-blocks}
\end{equation}
Rows index equations and columns index variables. We first establish that
these formal derivatives are the derivatives of the evaluation map.

\begin{lemma}[Derivative of polynomial evaluation]
\label{lem:actual-polynomial-derivative}
For $P\in\mathbb K[X_1,\ldots,X_n]$, polynomial evaluation is analytic at
every point. At $a\in\mathbb K^n$ its strict Fr\'echet derivative is
\[
 v\longmapsto\sum_{j=1}^n (\partial_jP)(a)v_j.
\]
For a finite system, its strict derivative is the continuous linear map
represented by the matrix of evaluated formal partial derivatives.
\end{lemma}
\begin{proof}
For the derivative statement, use polynomial induction on constants,
addition, and multiplication by a variable. Constants have zero derivative.
The sum rule agrees with the additivity of formal partial differentiation.
For $PX_j$, the product rule gives
\[
 d(PX_j)_a(v)=a_j\,dP_a(v)+P(a)v_j,
\]
which agrees with the evaluated formal product rule for every coordinate.
The finite product rule for maps to coordinate spaces then assembles the
system derivative. Each matrix defines a continuous linear map because the
coordinate spaces are finite-dimensional over the complete field.
Polynomial analyticity follows from the corresponding polynomial induction
in Mathlib and its finite-coordinate assembly.
\end{proof}
The derivative bridge and analytic statements are
\lean{Komorebi.PolynomialAnalyticEvaluation.hasStrictFDerivAt_evalSystem},
\lean{Komorebi.PolynomialAnalyticEvaluation.analyticAt_evalSystem}, and
\lean{Komorebi.PolynomialAnalyticEvaluation.analyticOnNhd_evalSystem}.

\begin{lemma}[The augmented derivative]
\label{lem:augmented-analytic-inverse}
If $J\in\operatorname{GL}_e(\mathbb K)$, the polynomial map
\[
 \Phi(Y,T)=(F(Y,T),T)
\]
has invertible strict derivative at $a$, with
\begin{equation}
 D\Phi(a)=\begin{pmatrix}J&K\\0&I_k\end{pmatrix},\qquad
 D\Phi(a)^{-1}=\begin{pmatrix}J^{-1}&-J^{-1}K\\0&I_k\end{pmatrix}.
 \label{eq:analytic-block-inverse}
\end{equation}
The local inverse supplied by this derivative is analytic near $\Phi(a)$.
\end{lemma}
\begin{proof}
The lower coordinates of $\Phi$ are the parameter variables. Differentiating
them gives the blocks $0$ and $I_k$. Lemma~\ref{lem:actual-polynomial-derivative}
identifies the remaining blocks with~\eqref{eq:regularity-blocks}.
Multiplication of the two matrices in~\eqref{eq:analytic-block-inverse},
in both orders, gives the identity. The resulting linear equivalence is
continuous, as is its inverse. Apply the strict inverse-function theorem to
$\Phi$. The analytic inverse theorem applies to the same open partial
homeomorphism, because $\Phi$ is analytic and its derivative is that
continuous linear equivalence. It supplies a convergent local power-series
representation of this inverse.
\end{proof}
The augmented Jacobian is identified in
\lean{Komorebi.PolynomialJointSystem.jacobian_augmented}.
The continuous inverse block and its restriction are
\lean{Komorebi.PolynomialAnalyticBlocks.upperEquiv} and
\lean{Komorebi.PolynomialAnalyticBlocks.inverse_solved_zero}.
Analyticity of the same constructed inverse is
\lean{Komorebi.PolynomialAnalyticInverse.analyticAt_inverse}.

\begin{theorem}[Analytic local root graph]
\label{thm:analytic-local-graph}
Suppose $F(a)=0$ and $J\in\operatorname{GL}_e(\mathbb K)$. There are open
neighborhoods $t_0\in U\subseteq\mathbb K^k$ and
$y_0\in V\subseteq\mathbb K^e$, and a map
$\psi:\mathbb K^k\to\mathbb K^e$, such that
\begin{gather}
 \psi(t_0)=y_0,\qquad \psi(U)\subseteq V,\qquad
 \psi\text{ is analytic at every point of }U,\label{eq:analytic-psi}\\
 F(y,t)=0\quad\Longleftrightarrow\quad y=\psi(t)
 \qquad(t\in U,\ y\in V).\label{eq:analytic-root-graph}
\end{gather}
At the base parameter, the strict Fr\'echet derivative is
\begin{equation}
 \boxed{D\psi(t_0)=-J^{-1}K.}
 \label{eq:implicit-derivative}
\end{equation}
In particular, $\psi$ is differentiable on $U$.
\end{theorem}
\begin{proof}
Let $g$ be the analytic local inverse of $\Phi$ supplied by
Lemma~\ref{lem:augmented-analytic-inverse}. Since $F(a)=0$,
$\Phi(a)=(0,t_0)$ and $g(0,t_0)=a$. Define
\[
 \psi(t)=\pi_Yg(0,t),
\]
where $\pi_Y$ projects to the solved coordinates. This gives the base value
and analyticity at $t_0$. The chain rule and
\eqref{eq:analytic-block-inverse} give
\[
 D\psi(t_0)v
 =\pi_Y\begin{pmatrix}J^{-1}&-J^{-1}K\\0&I_k\end{pmatrix}
       \binom{0}{v}
 =-J^{-1}Kv.
\]

To obtain the root equivalence on a product neighborhood, we choose the
neighborhoods as follows. The source of the open partial homeomorphism defined by
$\Phi$ is open around $a$. Its inverse is $g$. Choose open sets $V$ and $U_0$
containing $y_0$ and $t_0$ such that
$V\times U_0$ lies in that source. Restrict $U_0$ further by the preimage of
an open analytic inverse neighborhood under $t\mapsto(0,t)$, the open locus
where this same $\psi$ is analytic, and the interior of $\psi^{-1}(V)$.
All of these sets contain $t_0$. For the last set, use the continuity of
$\psi$ at $t_0$ and $\psi(t_0)=y_0\in V$. Their intersection is an open
neighborhood $U$ with $\psi(U)\subseteq V$.

If $t\in U$, $y\in V$ and $F(y,t)=0$, then
$\Phi(y,t)=(0,t)$. The left inverse identity forces
$g(0,t)=(y,t)$, hence $y=\psi(t)$. Conversely, the right inverse identity
$\Phi(g(0,t))=(0,t)$ says both that its solved equations vanish and that its
parameter coordinates are exactly $t$. Thus
$g(0,t)=(\psi(t),t)$ and $F(\psi(t),t)=0$.
Analyticity at every point of $U$ implies differentiability there.
\end{proof}
This is
\lean{Komorebi.PolynomialAnalyticImplicit.exists_analytic_implicit}.
Differentiability is also included explicitly in the statement contract
\lean{Komorebi.PolynomialRegularityContracts.actual_padic_analytic_root_graph}.

\begin{theorem}[Regularity of a supplied root selection]
\label{thm:selection-regularity}
Under the hypotheses of Theorem~\ref{thm:analytic-local-graph}, let
$h:\mathbb K^k\to\mathbb K^e$ be continuous at $t_0$, with
$h(t_0)=y_0$. Suppose $F(h(t),t)=0$ for every $t$ in some neighborhood of
$t_0$. Then $h$ and $\psi$ agree in a neighborhood of $t_0$. Consequently
$h$ is analytic there as a germ and has strict derivative
$Dh(t_0)=-J^{-1}K$.
\end{theorem}
\begin{proof}
Continuity at $t_0$ places $h(t)$ inside $V$ for all sufficiently close
parameters. Intersect this neighborhood with $U$ and the neighborhood on
which the root equation holds. Equation~\eqref{eq:analytic-root-graph} gives
$h(t)=\psi(t)$ throughout the intersection. Analyticity and strict
differentiability at a point are unchanged by agreement on a neighborhood.
\end{proof}
The formal result is
\lean{Komorebi.PolynomialAnalyticSelection.continuous_selection_regularity}.
Its formulation derives $F(a)=0$ from the root equation near the base
parameter and the specified base value.

\begin{corollary}[Regularity throughout an open parameter set]
\label{cor:selection-open-regularity}
Let $U\subseteq\mathbb K^k$ be open, and let $h$ be continuous on $U$.
Suppose that for every $t\in U$,
\[
 F(h(t),t)=0,
 \qquad J_t=D_YF(h(t),t)\in\operatorname{GL}_e(\mathbb K).
\]
Then $h$ is analytic at every point of $U$ and its strict derivative is
\[
 Dh(t)=-J_t^{-1}D_TF(h(t),t)\qquad(t\in U).
\]
\end{corollary}
\begin{proof}
At each $t\in U$, openness turns continuity on $U$ into continuity at $t$
and makes the root equation hold on a neighborhood of $t$. Apply
Theorem~\ref{thm:selection-regularity} at $(h(t),t)$.
\end{proof}
This is
\lean{Komorebi.PolynomialAnalyticSelection.continuous_selection_analyticOnNhd}.
The pointwise invertibility assumption is part of this corollary.

\begin{corollary}[Integral coefficients and centers]
\label{cor:integral-analytic-graph}
Let $p$ be any prime. Suppose $F\in\Zp[Y,T]^e$, $a\in\Zp^{e+k}$,
and $F(a)=0$. Let $\iota:\Zp\hookrightarrow\Qp$ be the inclusion.
If $\iota(D_YF(a))\in\operatorname{GL}_e(\Qp)$, the polynomial system
obtained by including the original coefficients in $\Qp$ has the analytic
local root graph of Theorem~\ref{thm:analytic-local-graph} at $\iota(a)$.
The strict derivative at the base is
\[
 -\iota(D_YF(a))^{-1}\iota(D_TF(a)).
\]
The open neighborhoods and the output map are in the ambient spaces over
$\Qp$.
\end{corollary}
\begin{proof}
Coefficient inclusion commutes with evaluation and formal partial
differentiation. Consequently both Jacobian blocks after inclusion are the
entrywise inclusions of the original blocks, and the original exact root
remains an exact root. Apply Theorem~\ref{thm:analytic-local-graph} to the system with included
coefficients.
\end{proof}
This is
\lean{Komorebi.PolynomialRegularityIntegral.exists_analytic_integral_implicit}.
It requires no invertibility modulo $p$. In particular, it includes solved
Jacobians singular modulo $p$ but invertible over the fraction field.

\begin{example}[The derivative sign and a nonlinear control]
For $F(Y,T)=Y+T$, the exact branch is $\psi(T)=-T$ and $J=K=1$.
The derivative is $-1$. Uniqueness of the derivative rules out $+1$ over
$\Qp$, including $p=2$. The nonlinear polynomial $F(Y,T)=Y+Y^2+T$
has $F(0,0)=0$ and $J=K=1$ there. Applying the formal theorem to this polynomial constructs an analytic
local root graph with derivative $-1$. Thus the example checks the existence
conclusion as well as the derivative formula.
The declarations and the deliberately rejected positive-sign statement are
in the regularity control and statement-contract files.
\end{example}

The results cover arbitrary finite $e$ and $k$, including zero. The local
inverse machinery explicitly handles a zero-dimensional domain. The analytic theorem does not assign a uniform numerical radius to its
neighborhoods. The earlier
weighted lifting and joint deformation theorems retain their quantitative
radii and error bounds, while the present statements establish regularity
of the local solution germ. A direct Lean identification with each
historical subtype-valued chart is an additional interface task. General
persistence at a fraction-field singular solved block is a separate problem.
